\section{实验设计与设置}
\label{sec:experiment}

本章从数据集、评价指标、实验环境及超参数配置四个方面说明实验的设计。

\subsection{数据集}
\label{sec:dataset}

实验数据来自 HuggingFace 开源数据集 \texttt{lansinuote/ChnSentiCorp}，包含酒店、笔记本电脑、书籍等场景的中文用户评论，每条标注为正向（1）或负向（0）。这是一个在中文情感分析中常用的基准数据集。

原始数据以 Parquet 格式存储，划分为 \texttt{train}（9,600 条）、\texttt{validation}（1,200 条）和 \texttt{test}（1,200 条）三个子集。实验中，SVM 将训练集与验证集合二为一（共 10,800 条）进行训练，不设独立验证集；Bi-LSTM、BERT 与 Qwen-LoRA 使用训练集（9,600 条）训练，验证集（1,200 条）用于超参数监控与最优模型选择。所有模型在测试集（1,200 条）上评估。具体划分如表~\ref{tab:dataset}~所示。

\begin{table}[htbp]
    \centering
    \caption{ChnSentiCorp 数据集划分}
    \label{tab:dataset}
    \begin{tabular}{ccccc}
        \toprule
        子集 & SVM & Bi-LSTM & BERT/Qwen-LoRA \\
        \midrule
        训练集 & 10,800 & 9,600 & 9,600\\
        验证集 & ---    & 1,200 & 1,200\\
        测试集 & 1,200  & 1,200 & 1,200\\
        \midrule
        总计   & 12,000 & 12,000 & 12,000\\
        \bottomrule
    \end{tabular}
\end{table}

在预处理阶段，所有文本经过中文分词（SVM 与 Bi-LSTM 使用 Jieba 分词，BERT 与 Qwen 使用各自的 Tokenizer）。

\subsection{评价指标}
\label{sec:metrics}

情感分类本质上是一个二分类问题。本文采用准确率（Accuracy）、精确率（Precision）、召回率（Recall）和 F1-Score 四项指标进行评估：

\begin{equation}
    \text{Accuracy} = \frac{\text{TP} + \text{TN}}{\text{TP} + \text{TN} + \text{FP} + \text{FN}}
\end{equation}

\begin{equation}
    \text{Precision} = \frac{\text{TP}}{\text{TP} + \text{FP}}
    \label{eq:precision}
\end{equation}

\begin{equation}
    \text{Recall} = \frac{\text{TP}}{\text{TP} + \text{FN}}
    \label{eq:recall}
\end{equation}

\begin{equation}
    \text{F1} = \frac{2 \times \text{Precision} \times \text{Recall}}{\text{Precision} + \text{Recall}}
\end{equation}

其中，TP 为真正例，TN 为真负例，FP 为假正例，FN 为假负例。准确率反映整体正确率；精确率衡量模型判为正类时的可靠性；召回率衡量对正类样本的覆盖率；F1-Score 是精确率与召回率的调和平均，用于综合评价分类质量。此外，实验同时记录了各模型的训练耗时和推理耗时。

\subsection{实验环境}
\label{sec:environment}

所有模型的训练和测试在统一的硬件与软件环境下完成，配置如表~\ref{tab:env}~。

\begin{table}[htbp]
    \centering
    \caption{实验软硬件环境}
    \label{tab:env}
    \begin{tabular}{ll}
        \toprule
        组件 & 配置 \\
        \midrule
        GPU          & NVIDIA RTX PRO 6000 Blackwell Server Edition (96 GB VRAM) \\
        CPU          & Intel(R) Xeon(R) Platinum 8470Q \\
        内存         & 110 GB \\
        操作系统     & Ubuntu 22.04 \\
        Python       & 3.12 \\
        PyTorch      & 2.11.0 + CUDA 12.8 \\
        Transformers & 5.8.0 \\
        Scikit-learn & 1.8.0 \\
        \bottomrule
    \end{tabular}
\end{table}

\subsection{模型配置与超参数}
\label{sec:hyperparams}

为保证实验公平，各模型的超参数在预实验调优后确定，全程保持不变。

\subsubsection{SVM + TF-IDF 基线模型}

使用 TF-IDF 向量化器对分词后的文本做特征提取，最大特征数设为 10,000。分类器选用 LinearSVC，惩罚系数 $C = 1.0$，最大迭代次数 2,000，其余参数为 Scikit-learn 默认值。

\subsubsection{Bi-LSTM 模型}

Bi-LSTM 从零学习词嵌入。词表大小 15,000（含 \texttt{<PAD>} 和 \texttt{<UNK>} 两个特殊标记），嵌入维度与隐藏层维度均为 128。LSTM 采用单层双向结构，在分类层前施加 Dropout（比率 0.5）。优化器为 Adam，初始学习率 $1 \times 10^{-3}$，批次大小 64。训练引入早停机制：最大训练轮数 50，验证集损失连续 5 轮不下降则终止训练，取验证集准确率最高的检查点作为最终模型。

\subsubsection{BERT 微调模型}

模型采用 \texttt{bert-base-chinese} 预训练权重：12 层 Transformer 编码器，隐藏维度 768，总参数量约 1.1 亿。微调使用 HuggingFace Trainer 框架，学习率 $2 \times 10^{-5}$，批次大小 32，最大序列长度 128。训练引入早停机制（EarlyStoppingCallback）：最大训练轮数 20，耐心值 3，验证集指标连续 3 轮不提升则终止。预热步数比例（Warmup Step Ratio）为 0.1，权重衰减为 0.01。每轮评估并保存检查点，训练结束后加载验证集最优检查点。

\subsubsection{Qwen-7B 零样本推理}

零样本实验使用 \texttt{Qwen2.5-7B} 预训练权重，以 BF16 精度加载，通过 \texttt{device\_map="auto"} 自动分配至 GPU。该模型拥有 70 亿参数、32 层 Transformer 解码器、32 个注意力头。推理时采用 Chat Template：系统消息为"情感分析：好评输出1，差评输出0。"，用户消息为待分析文本，随后通过规则提取标签（检测输出中是否含'1'或'正'判为正向，'0'或'负'判为负向）。最大生成 Token 数 5，关闭随机采样（\texttt{do\_sample=False}）。此实验不更新模型权重。

\subsubsection{Qwen-7B LoRA 微调}

LoRA 微调同样基于 \texttt{Qwen2.5-7B} 预训练权重，以 BF16 精度加载。模型头部替换为二分类的 \texttt{AutoModelForSequenceClassification} 结构，设 \texttt{pad\_token\_id} 为 \texttt{eos\_token\_id}。

LoRA 配置：秩 $r = 16$，缩放因子 $\alpha = 32$，Dropout 比率 0.05，目标模块为注意力层的 Query 投影（\texttt{q\_proj}）和 Value 投影（\texttt{v\_proj}）。训练时冻结原始模型全部权重，仅优化 LoRA 新增参数。

训练参数：批次大小 16，梯度累积步数 2（等效批次 32），学习率 $1 \times 10^{-4}$，预热步数比例 0.05。引入早停机制：最大训练轮数 20，耐心值 3。启用 BF16 混合精度训练，每轮评估并保存检查点，结束加载最优检查点。
